% from ega4-2-fr.tex
\subsection{Application aux changements de base dans les préschémas algébriques.}

\begin{proposition}[6.7.1]
\label{IV.2.s6.7.1-fr}
Soient $k$ un corps, $X$ un $k$-préschéma localement noethérien,
$\mathscr F$ un $\mathcal O_X$-Module cohérent. Soit $k'$ une extension
de $k$; posons $X'=X\otimes_k k'$, $\mathscr F'=\mathscr F\otimes_k k'$
et soit $p:X'\to X$ la projection canonique. On suppose, soit que
$X$ est localement de type fini sur $k$, soit que $k'$ est une extension
de type fini de $k$, de sorte que dans les deux cas $X'$ est localement
noethérien. Soient $x'$ un point de $X'$, $x=p(x')$. Alors:
\begin{enumerate}
\item[\textup{(i)}] On a $\operatorname{coprof}(\mathscr F_x)=
\operatorname{coprof}(\mathscr F'_{x'})$; en particulier, pour que
$\mathscr F_x$ soit un $\mathcal O_x$-module de Cohen-Macaulay, il
faut et il suffit que $\mathscr F'_{x'}$ soit un
$\mathcal O_{x'}$-module de Cohen-Macaulay.
\item[\textup{(ii)}] Pour que $\mathscr F$ vérifie la propriété $(S_n)$ au point
$x$, il faut et il suffit que $\mathscr F'$ vérifie $(S_n)$ au point
$x'$.
\end{enumerate}
\end{proposition}

\oldpage[IV-2]{146}

\begin{proof}
On sait que $p$ est fidèlement plat (2.2.13); les deux assertions sont
donc conséquences respectives de (6.3.2) et (6.4.2), pourvu que l'on
prouve que $p^{-1}(x)=\operatorname{Spec}(\boldsymbol{k}(x)\otimes_k k')$ est un
\emph{préschéma de Cohen-Macaulay}; comme l'un des deux corps $\boldsymbol{k}(x)$, $k'$ est
une extension de type fini de $k$ par hypothèse on est ramené à prouver
le
\end{proof}

\begin{lemma}[6.7.1.1]
\label{IV.2.s6.7.1.1-fr}
Soient $K$, $L$ deux extensions d'un corps $k$, dont l'une est de type
fini (de sorte que l'anneau $K\otimes_k L$ est noethérien). Alors
$K\otimes_k L$ est un anneau de Cohen-Macaulay.
\end{lemma}

\begin{proof}
Supposons par exemple que $L$ soit une extension de type fini de $k$,
de sorte que $L$ est extension finie d'une extension pure $k(\mathbf t)$
de $k$ ($\mathbf t$ étant un système $(t_i)_{1\leqslant i\leqslant m}$
d'indéterminées). Si on pose $A=K\otimes_k k(\mathbf t)$,
$B=K\otimes_k L$, $B$ est un $A$-module plat $(\mathbf{0}_{\mathrm I},6.2.1)$ et de type fini;
le morphisme $h:\operatorname{Spec}(B)\to\operatorname{Spec}(A)$ est
donc fini, et comme tout préschéma artinien est de Cohen-Macaulay, pour
prouver que $B$ est un anneau de Cohen-Macaulay, il suffit de prouver que
$A$ est un anneau de Cohen-Macaulay (6.3.5). Or, si $S$ est l'ensemble
des éléments $\neq 0$ de $k[\mathbf t]$, on a $A=S^{-1}A'$, où
$A'=K\otimes_k k[\mathbf t]=K[\mathbf t]$; en vertu de
(\textbf{0}, 16.5.13), il suffit de prouver que $A'$ est un anneau de
Cohen-Macaulay; mais cela résulte de ce que $A'$ est \emph{régulier} (\textbf{0}, 17.3.7).
\end{proof}

\begin{corollary}[6.7.2]
\EGAReaderTarget{ega4.c04.cor.6-7-2}{IV.2.s6.7.2-fr}
Pour que $\mathcal O_x$ soit un anneau de Cohen-Macaulay \textup{(resp.} pour
que $X$ vérifie $(S_{\boldsymbol{n}})$ au point $x$), il faut et il suffit
que $\mathcal O'_{x'}$ soit un anneau de Cohen-Macaulay \textup{(resp.} que $X'$
vérifie $(S_{\boldsymbol{n}})$ au point $x'$).
\end{corollary}

\begin{corollary}[6.7.3]
\label{IV.2.s6.7.3-fr}
Soient $X$, $Y$ deux $k$-préschémas localement noethériens, l'un au
moins étant localement de type fini sur $k$. Soit $\mathscr F$ \textup{(resp.}
$\mathscr G$) un $\mathcal O_X$-Module cohérent \textup{(resp.} un
$\mathcal O_Y$-Module cohérent). Si $\mathscr F$ et $\mathscr G$
vérifient la propriété $(S_n)$, il en est de même de
$\mathscr F\otimes_k\mathscr G$.
\end{corollary}

\begin{proof}
L'hypothèse entraîne que $X\times_kY$ est localement noethérien; notons
$p:X\times_kY\to X$ et $q:X\times_{\boldsymbol k}Y\to Y$ les projections canoniques,
qui sont des morphismes plats. Supposons par exemple que $X$ soit
localement de type fini sur $k$; on peut écrire
$\mathscr F\otimes_k\mathscr G=p^*(\mathscr F)\otimes_k\mathscr G$
et puisque $\mathscr F$ est plat relativement au morphisme structural
$X\to\operatorname{Spec}(k)$, $p^*(\mathscr F)$ est $q$-plat (2.1.4).
Appliquons le critère (6.4.1, (ii)) au morphisme $q$; il suffit de voir
que pour tout $y\in Y$, $p^*(\mathscr F)\otimes_{\mathcal O_Y}\boldsymbol{k}(y)$
a la propriété $(S_n)$; mais
$p^*(\mathscr F)\otimes_{\mathcal O_Y}\boldsymbol{k}(y)=\mathscr F\otimes_k \boldsymbol{k}(y)$,
et comme $X$ est localement de type fini sur $k$, la conclusion résulte
de (6.7.1, (ii)).
\end{proof}

\begin{proposition}[6.7.4]
\label{IV.2.s6.7.4-fr}
Pour un préschéma localement noethérien $Z$, et un point $z\in Z$,
soit $\boldsymbol{P}(Z,z)$ l'une des propriétés suivantes:
\begin{enumerate}
\item[\textup{a)}] $Z$ est régulier au point $z$.
\item[\textup{b)}] $Z$ vérifie $(R_n)$ au point $z$.
\item[\textup{c)}] $Z$ est normal au point $z$.
\item[\textup{d)}] $Z$ est réduit au point $z$.
\end{enumerate}
Cela étant, sous les hypothèses et avec les notations de \textup{(6.7.1)}, si
$\boldsymbol{P}(X',x')$ est vraie, il en est de même de $\boldsymbol{P}(X,x)$. La réciproque est
vraie si $k'$ est une extension séparable de $k$.
\end{proposition}

\begin{proof}
Le cas où $\boldsymbol{P}$ est la propriété \emph{d)} n'a été mis que pour mémoire, ayant
déjà été traité (2.1.13 et 4.6.1). Le même raisonnement qu'au début
de (6.7.1) montre que la première assertion résulte
respectivement de (6.5.2, (i)), (6.5.3, (i)) et (6.5.4, (i)); de même,
la seconde assertion résultera de (6.5.2, (ii)), (6.5.3, (ii)) et
(6.5.4, (ii))
